% ============================================================================
% Section 4.6 — Rewriting Expressions to Solve Problems
% CCSS 7.EE.A.2
% ============================================================================
\section{Rewriting Expressions to Solve Problems}

\begin{stepGoal}
\begin{itemize}
    \item Rewrite expressions in different forms to reveal useful information
    \item Connect equivalent expressions to real-world meaning
\end{itemize}
\end{stepGoal}

\begin{stepVocabBox}
    \vocabItem{Equivalent Expressions}{Two expressions that have the \textbf{same value} for every number you plug in.}
    \vocabItem{Rewrite}{Change the form of an expression without changing its value.}
\end{stepVocabBox}

\begin{stepsCard}[How to Rewrite Expressions to Solve Problems]
    \stepItem{Write the expression that \textbf{matches the situation} (even if it looks long).}
    \stepItem{Use the distributive property, combining like terms, or factoring to \textbf{simplify} to a shorter form.}
    \stepItem{Read the new form to find the \textbf{hidden meaning} (e.g.\ a single multiplier, a percent, a total).}
\end{stepsCard}

% --- Worked Example 1 ---
\begin{stepExample}{A store marks up a price $p$ by $8\%$. Write two equivalent expressions for the new price.}
    \stepShow{1} Start price is $p$. The markup is $8\%$ of $p$: new price $= p + 0.08p$.

    \stepShow{2} Factor out $p$: $p(1 + 0.08) = 1.08p$.

    \stepShow{3} The form $1.08p$ shows the new price is $108\%$ of the original --- one multiplication instead of two steps.
    \stepResult{$p + 0.08p = 1.08p$}
\end{stepExample}

% --- Worked Example 2 ---
\begin{stepExample}{Simplify $3(x + 4) - 2x$. What does the result tell you?}
    \stepShow{1} The expression is $3(x + 4) - 2x$.

    \stepShow{2} Expand: $3x + 12 - 2x$. Combine like terms: $x + 12$.

    \stepShow{3} The simpler form $x + 12$ shows the result is always $12$ more than $x$.
    \stepResult{$x + 12$}
\end{stepExample}

\mascotStepTip{``Increase by $5\%$'' means multiply by $1.05$. ``Decrease by $20\%$'' means multiply by $0.80$. One number does it all!}

% ============================================================================
% PRACTICE
% ============================================================================
\newpage
\begin{stepPractice}{Rewriting Expressions Practice}
    \resetProblems

    \stepPracticeHeader[step1Color]{Write the Starting Expression}
    \prob A price $p$ is discounted $15\%$. Write an expression for the discount and the sale price. \answerBlank[3cm]
    \answer{Discount: $0.15p$; sale price: $p - 0.15p$}

    \stepPracticeHeader[step2Color]{Simplify to a Single Term}
    \begin{multicols}{2}
    \prob Rewrite $x + 0.08x$. \answerBlank[2.5cm]
    \answer{$1.08x$}
    \prob Rewrite $p - 0.15p$. \answerBlank[2.5cm]
    \answer{$0.85p$}
    \end{multicols}

    \stepPracticeHeader[step3Color]{Interpret the Meaning}
    \prob Rewrite $3(x + 4) - 2x$ in simplest form. \answerBlank[2.5cm]
    \answer{$x + 12$}

    \prob A worker earns $\$h$ per hour for $40$ hours plus a $\$50$ bonus. Write a simplified expression for weekly pay. \answerBlank[3cm]
    \answer{$40h + 50$}

    \wordProblem{Show that ``increase by $25\%$ then decrease by $20\%$'' returns the original price $p$.}{~}
    \answer{$1.25p \times 0.80 = 1.00p = p$}
\end{stepPractice}
